\begin{abstract}

% Recent works show that near-oracle top-k KV selection for the attention operator can match the generation quality of full-cache attention while sharply reducing computation and bandwidth. Yet because computing the oracle itself requires scoring all KV entries, prior work has resorted to posterior heuristics—selectors conditioned on observed attention or proxy scores.
% However, such conditioning introduces posterior bias: it tends to distort the true token importance and miss salient tokens, impairing long-range reasoning.

A core bottleneck in large language model (LLM) inference is the cost of attending over the ever-growing key–value (KV) cache. Although near-oracle top-$k$ KV selection can preserve the quality of dense attention while sharply reducing computation and bandwidth, existing sparse methods generally rely on posterior heuristics—selectors conditioned on observed attention or proxy scores. Such conditioning introduces posterior bias: it tends to distort the true token importance and miss salient tokens, thereby impairing long-range reasoning.
To tackle this problem, we propose Pre-hoc Sparsity (PrHS), which selects KV entries before attention scoring and provides explicit accuracy control. Let the attention mass of discarded entries be $\delta$ (the dropped mass). Through a marginal-to-mutual-information analysis, we derive an upper bound on the mutual-information loss that depends only on $\delta$. This relation both explains the failure modes of posterior heuristics and enables verifiable guarantees by controlling $\delta$ in advance. Within PrHS, we instantiate three orthogonal pre-hoc selectors along the axes of time, depth, and layer: (i) temporal query sharing, justified by a Lipschitz bound on attention-centroid drift so that a modest dilation of the previously shared set covers the next step; (ii) recency-window retention, leveraging distance-decay priors to guarantee retained attention mass; and (iii) cross-layer KV sharing, exploiting high inter-layer similarity to eliminate redundant retrieval and computation. Extensive experiments on LLaMA and Mistral families validate the superiority of PrHS.  On LLaMA and Mistral families across GSM8K and CoQA, the approach reduces retrieval overhead by ${>}90\%$, achieving $3\times$ higher retrieval sparsity than the state-of-the-art HShare at matched or better accuracy. It also incurs ${<}1\%$ average degradation on LongBench, lowers attention FLOPs by ${\sim}15\%$ versus prior sparse baselines, and yields a $9.9\times$ speedup in attention-operator latency and $2.8\times$ higher throughput on NVIDIA A100-80GB GPUs than the original dense baseline.
\end{abstract}


% In long-context inference, large language models (LLMs) suffer from key-value (KV) cache growth that scales linearly with sequence length. KV compression leverages attention sparsity by retaining only critical KV entries. However, since the optimal top-$k$ set is intractable to identify, forcing heuristic approximations is a necessity. From a new marginal entropy perspective, we reinterpret existing approximations as prefixing simplifications grounded from posterior observations, which largely distort the true token importance. To eliminate this posterior bias, we devise Pre-hoc Sparsity (PrHS), a principled framework that selects KV entries according to the pre-hoc bounded marginal entropy and provides a theoretical accuracy guarantee via a mutual information upper bound. 